\chapter{Esquema de campaña}
\label{anx:esquema}

Las campañas de todos los patrones del Pilar~A se emiten como hallazgos \textbf{OCSF 1.5.0
«Detection Finding»} (\texttt{class\_uid} 2004, \texttt{type\_uid} 200401) con una extensión
propia, \texttt{tfg}. Se publican en el \textit{topic} \texttt{det.campaigns}, sin caducidad,
y se guardan en la tabla \texttt{tfg.campaigns} del almacén analítico.

\section{Campos del estándar}

\begin{table}[H]
\centering
\small
\begin{tabular}{|l|p{8.2cm}|}
\hline
\textbf{Campo} & \textbf{Contenido} \\
\hline
\texttt{start\_time}, \texttt{end\_time}, \texttt{duration} & Inicio, fin y duración de la campaña, en milisegundos \\
\texttt{time} & Igual a \texttt{end\_time}, para que reprocesar produzca el mismo documento \\
\texttt{severity\_id} & Máxima gravedad de las reglas que aportan evidencia \\
\texttt{confidence\_score} & Suma de los factores de confianza, de 0 a 100 \\
\texttt{resources[]} & Equipos implicados \\
\texttt{finding\_info.uid} & Identidad de la campaña \\
\texttt{finding\_info.analytic} & Patrón que la reconoce (P1, P2, P3) \\
\texttt{finding\_info.attacks[]} & Técnicas observadas en la evidencia, sin táctica asociada \\
\texttt{finding\_info.related\_events[]} & Alertas y eventos de evidencia, una vez cada uno \\
\hline
\end{tabular}
\caption{Campos del estándar OCSF empleados.}
\end{table}

\section{Extensión propia}

\begin{table}[H]
\centering
\small
\begin{tabular}{|l|p{9cm}|}
\hline
\textbf{Campo} & \textbf{Contenido} \\
\hline
\texttt{tfg.pattern\_id}, \texttt{tfg.dataset\_id} & Patrón y conjunto de datos de origen \\
\texttt{tfg.devices[]} & Equipos con su papel: \textit{host} en P1 y P3, origen y destino en P2 \\
\texttt{tfg.tactic\_names[]} & Tácticas de la evidencia, por nombre \\
\texttt{tfg.stages[]} & Etapas ordenadas, cada una con su intervalo, sus técnicas y sus evidencias \\
\texttt{tfg.confidence\_factors} & Factores presentes y su peso; su suma es la confianza \\
\texttt{tfg.metrics} & Métricas propias del patrón (aristas, transiciones, alertas en la ventana) \\
\hline
\end{tabular}
\caption{Campos de la extensión \texttt{tfg}.}
\end{table}

\section{Reglas de construcción}

\begin{itemize}
    \item \textbf{Identidad}: \texttt{md5(dataset | patrón | clave | inicio en ms)}, donde la
    clave es el equipo en P1 y P3 y el par origen--destino en P2. Es determinista.
    \item \textbf{Orden de las etapas}: por el segundo de la primera evidencia de cada etapa y,
    dentro del mismo segundo, por el orden lógico del patrón.
    \item \textbf{Una evidencia puede probar dos etapas} (el acceso remoto a LSASS mediante
    WinRM es a la vez ejecución remota y acceso a credenciales). Aparece en las dos etapas,
    pero una sola vez en \texttt{related\_events}.
    \item \textbf{Validación}: el proceso de volcado comprueba la clase, la coherencia de los
    tiempos, que la confianza sea la suma de sus factores, que las etapas estén numeradas de 1
    a $n$ con su evidencia y que no haya evidencias repetidas.
\end{itemize}

\section{Ejemplo}

El Listado~\ref{lst:campana-p2} muestra, abreviada, la campaña de movimiento lateral del
primer día de APT29.

\begin{lstlisting}[caption={Campaña P2 de APT29, abreviada.},label={lst:campana-p2}]
{
  "class_uid": 2004, "type_uid": 200401, "activity_id": 1,
  "start_time": 1588388963801, "end_time": 1588389479493,
  "duration": 515692, "severity_id": 4, "confidence_score": 100,
  "message": "Movimiento lateral SCRANTON -> NASHUA",
  "finding_info": {
    "uid": "07c1426c17179d180b24073e097edfdc",
    "analytic": {"uid": "P2", "name": "P2_lateral_movement"},
    "attacks": [{"technique": {"uid": "T1021"},
                 "sub_technique": {"uid": "T1021.006"}}],
    "related_events_count": 8
  },
  "tfg": {
    "pattern_id": "P2",
    "devices": [{"hostname": "SCRANTON", "role": "origen"},
                {"hostname": "NASHUA", "role": "destino"}],
    "stages": [{"order": 1, "name": "conexion_lateral", "count": 17}],
    "confidence_factors": {"arista": 0.30, "bilateral": 0.15,
      "proceso_interactivo": 0.10, "ejecucion_remota": 0.20,
      "servicio": 0.15, "credenciales": 0.10},
    "metrics": {"aristas": 17, "transiciones": 13,
      "alertas_absorbidas": 115, "retraso_max_ms": 1004}
  }
}
\end{lstlisting}
